\section{Descripción en HDL y Hardware}

\subsection{Código Verilog}
La descripción del hardware se realizó utilizando el lenguaje Verilog. El módulo principal, denominado \texttt{monta}, implementa la máquina de estados descrita anteriormente.

\subsubsection{Código}

\begin{lstlisting}
module monta(
	input clk,
	input reset,
	input P1, P2, P3,
	input fc1, fc2, fc3,
	output en,
	output reg [1:0] S,
	output reg [6:0] D 
);

assign en = 0;

reg [2:0] state, nextState;

parameter piso1 = 3'b000, piso2 = 3'b010, piso3 = 3'b100;
parameter sube2 = 3'b001, baja1 = 3'b011;
parameter sube3 = 3'b101, baja2 = 3'b110;

always @(posedge clk or posedge reset) begin
	if (reset) state <= piso1;
	else state <= nextState;
end

always @(*) begin
    nextstate = state;
    S = 2'b00;
    D = 7'b0000000;

    case (state)
        piso1: begin
            D = 7'b0110000;
            if (P2)      begin nextstate = sube2; S = 2'b01; D = 7'b0000000; end
            else if (P3) begin nextstate = sube3; S = 2'b01; D = 7'b0000000; end
        end
        sube2: begin
            S = 2'b01; 
            if (fc2) begin nextstate = piso2; S = 2'b00; D = 7'b1101101; end
        end
        sube3: begin
            S = 2'b01; 
            if (fc3) begin nextstate = piso3; S = 2'b00; D = 7'b1111001; end
        end
        piso2: begin
            D = 7'b1101101; 
            if (P1)      begin nextstate = baja1; S = 2'b10; D = 7'b0000000; end
            else if (P3) begin nextstate = sube3; S = 2'b01; D = 7'b0000000; end
        end
        baja1: begin
            S = 2'b10; 
            if (fc1) begin nextstate = piso1; S = 2'b00; D = 7'b0110000; end
        end
        piso3: begin
            D = 7'b1111001;
            if (P1)      begin nextstate = baja1; S = 2'b10; D = 7'b0000000; end
            else if (P2) begin nextstate = baja2; S = 2'b10; D = 7'b0000000; end
        end
        baja2: begin
            S = 2'b10; 
            if (fc2) begin nextstate = piso2; S = 2'b00; D = 7'b1101101; end
        end
        default: nextstate = piso1;
endcase
end
endmodule

\end{lstlisting}

\subsubsection{Estructura del Módulo}
El diseño se divide en dos bloques principales dentro del módulo:

\begin{enumerate}
    \item \textbf{Lógica Secuencial (Memoria de Estado):}
    Se utiliza un bloque \texttt{always @ (posedge clk, posedge reset)} para actualizar el estado del sistema.
    \begin{itemize}
        \item Si la señal de \texttt{reset} es alta, el sistema se fuerza asincrónicamente al estado \texttt{piso1} (estado inicial).
        \item En cada flanco positivo del reloj (\texttt{clk}), el registro \texttt{state} toma el valor de \texttt{nextState}.
    \end{itemize}

    \item \textbf{Lógica Combinacional (Estado Siguiente y Salida):}
    Se utiliza un bloque \texttt{always @(...)} sensible a todas las entradas y al estado actual. Mediante una sentencia \texttt{case}, se evalúa el estado presente para determinar:
    \begin{itemize}
        \item El valor de \texttt{nextState} según las entradas de los pulsadores ($P$) y finales de carrera ($fc$).
        \item Los valores de las salidas $S$ (motor) y $D$ (display).
    \end{itemize}
    Se definieron parámetros (ej. \texttt{parameter piso1 = 3'b000}) para hacer el código más legible y facilitar el mantenimiento.
\end{enumerate}

\subsection{Implementación en FPGA/CPLD}
El diseño fue sintetizado e implementado en el kit de desarrollo de la cátedra, basado en un CPLD de Xilinx. La asignación de pines (User Constraints) se realizó de acuerdo a la disposición física de los periféricos en la placa.

\imagen[Implementación en CPLD]{15 cm}{./imagenes/cpld}

\subsubsection{Asignación de Puertos}
A continuación se detalla el mapeo físico entre las señales del diseño Verilog y los pines del dispositivo CPLD:

\begin{table}[H]
\centering
\caption{Asignación de Pines en el CPLD}
\begin{tabular}{|c|c|l|}
\hline
\textbf{Señal Verilog} & \textbf{Pin CPLD} & \textbf{Descripción del Periférico} \\ \hline
\texttt{clk} & P6 & Señal de reloj (Frecuencia Variable) \\ \hline
\texttt{reset} & P38 & Switch para reinicio del sistema \\ \hline
\texttt{P1, P2, P3} & P34, P35, P36 & Switches para selección de piso \\ \hline
\texttt{fc1, fc2} & P27, P33 & Pulsadores (Lógica inversa: 0 presionado) \\ \hline
\texttt{fc3} & P37 & Switch (Simulación de final de carrera superior) \\ \hline
\texttt{S} & P14 & LED indicador de estado del motor \\ \hline
\texttt{en} & P28 & Habilitador del Display (Transistor) \\ \hline
\texttt{D[6:0]} & P18 - P26 & Segmentos del Display 7-Segmentos (a-g) \\ \hline
\end{tabular}
\label{tab:pines}
\end{table}

\imagen[Test bench]{15 cm}{./imagenes/test}

\subsubsection{Consideraciones de Hardware}
\begin{itemize}
    \item \textbf{Display de 7 Segmentos:} Los segmentos están conectados a los pines P18 a P26. El código envía los vectores de bits correspondientes para formar los números '1', '2' y '3'. El pin P28 controla el habilitador común del display.
    \item \textbf{Entradas de Usuario:} Se utilizaron switches (P34-P36) para simular las llamadas de los pisos, permitiendo mantener la señal activa.
    \item \textbf{Sensores de Posición:} Para los finales de carrera 1 y 2 se utilizaron los pulsadores de la placa (P27 y P33), los cuales operan con lógica negativa (envían un '0' al ser presionados). Para el final de carrera 3, se utilizó un switch (P37) por disponibilidad de hardware.
\end{itemize}
